\documentclass[11pt]{amsart}
\usepackage[margin=1in]{geometry}
\usepackage{amsmath,amssymb,array,booktabs,longtable}
\usepackage{xurl}
\usepackage[hidelinks,hypertexnames=false]{hyperref}
\newcommand{\Z}{\mathbb Z}
\newcommand{\N}{\mathbb N}
\newcommand{\GC}{\operatorname{GC}}
\newcommand{\record}[1]{\nolinkurl{#1}}
\newcommand{\claimid}[1]{\par\nopagebreak\noindent\textit{Machine-readable claim:} \record{#1}.\par}
\newtheorem{alignedtheorem}{Theorem}
\renewcommand{\thealignedtheorem}{1a}
\newtheorem{catalogcor}{Corollary}
\renewcommand{\thecatalogcor}{1b}
\newtheorem{proposition}{Proposition}
\newtheorem{lemma}{Lemma}
\theoremstyle{definition}\newtheorem{definition}{Definition}
\setlength{\emergencystretch}{3em}
\raggedbottom
\title[Affine fullerene families and eventual periodicity]{Eventual periodicity of linear generalized Stone--Wales paths in affine fullerene families with convex lattice-face decompositions}
\author{Tatsuhiko Takimoto}
\address{Independent researcher}
\date{Version 4.6, 20 October 2026}
\begin{document}
\begin{abstract}
We study literal linear generalized Stone--Wales paths in fullerene
duals. Our main result concerns an integer affine family admitting a
fixed decomposition into embedded convex lattice polygons with sides
in the six triangular-lattice directions. With fixed gluing and affine
corner coordinates, we give a uniform reduction of literal-path
existence to a fixed monadic second-order sentence over unary finite
words, using maximal face passages, their two-row outputs and all
vertex collisions. Eventual periodicity on every eventually valid
integer affine ray then follows from the known equivalence between
MSO and regular languages. This is a geometric definability reduction,
not a new general periodicity principle. We record
the resulting ray statements for eleven specified catalogue families,
alongside closed conditions on their stated domains. Existing Lean
results establish arithmetic periodicity of several family predicates
and determine the decisions on two example rays for every value of the ray parameter. For K3, both the graph classification and eventual periodicity on every
positive integer affine ray are now Lean-proved on the explicit quotient
specification, including parameter-one boundaries and without the general
prose bridge. The general eventual-periodicity statement is now proved
in Lean against a specification written by the coordinator before the
proof. Its correspondence to the geometric hypotheses is an explicit
prose lemma, not a Lean theorem. The proof was supplied by Claude
workers; no independent Lean proof by the research team is claimed.
The separate R004 prose derivation is retained. Prior owner confirmation of the general specification is retained as a
historical record, not a new verification procedure. Human proof peer review
and novelty remain open;
no executable geometric-to-MSO compiler or period algorithm is supplied.
\end{abstract}
\maketitle

\noindent\textbf{Editorial and evidence status.}
This English edition continues Japanese v3. The
main result is its convex-face Theorem 1a and catalogue Corollary 1b.
The wider fixed-oblique-direction argument is retained only as a proof
sketch in Section~\ref{sec:extension}, outside the abstract's claims.
The contribution is the uniform geometric reduction to a fixed MSO
sentence on unary words; the final periodicity implication is known.
Section~\ref{sec:novelty} separates that reduction from standard
MSO, parametric Presburger and fixed-boundary iteration results.
The independent verifier R004 derived the convex-face proof without
reading the research team's proof, construction or code. The statement
was known. Disclosures about O2 mapping metadata and D12 single-point
information are retained; strict blindness is not claimed. One AI
coordinator reviewed the route. Translation and packaging are not a
third mathematical derivation or human peer review. The exact
coordinator-written Theorem 1a specification now has a Lean proof
and a saved six-kernel Comparator acceptance. On 20 October 2026,
the owner reported reading version 2 with the Japanese guide and its
addendum and confirming that it expresses Theorem 1a. This checks
statement meaning, not the proof or prose correspondence lemma. Previous
authorship, affiliation and licensing decisions are retained; the
statement-meaning check is separate from publication decisions. On
20 October 2026 the owner discontinued owner statement-meaning checks as a
verification procedure; the earlier record remains factual. The K3
correspondence evidence instead comprises a specification written before
the proof, an independent AI statement reader and finite computational
comparisons. The first paper's
A3 classification is a separate work; its 816-class census is not
reproduced here. A new D12 classification is not incorporated.

\clearpage
\section{Definitions, signatures and valid rays}\label{sec:definitions}
\subsection{The literal path}
\begin{definition}\label{def:literal}
In the dual $G$ of a fullerene, a literal linear gSW path is a sequence
$p_0,\ldots,p_e$ with odd $e\ge3$, all vertices distinct, consecutive
and distance-two vertices adjacent, $\deg p_0=\deg p_e=5$, and
$\deg p_1=\deg p_{e-1}=6$.
\end{definition}
\claimid{def:literal}
This is Definition 1(a) of \cite{BBS}. Internal degree-five vertices
are allowed. Even vertex count is part of the definition, not a new
obstruction. Positive means existence of such a path; no chemical
barrier, radial operation, or operation permitting heptagons is inferred.
A fullerene is a simple three-connected cubic plane graph with only
pentagonal and hexagonal faces. Its dual is a spherical triangulation
with twelve degree-five vertices and all other degrees six.

When every three-cycle is facial and an edge has just the two common
neighbours supplied by its incident triangles, an oriented triangle
$(a,b,c)$ has a unique next vertex: the common neighbour of $b,c$
different from $a$. A simple path is a prefix of this forced sequence.
If $\rho$ is its first repeated-vertex index, a legal endpoint must
occur strictly before $\rho$, not at $\rho$.
\label{prop:forced}\claimid{prop:forced}
The conditional finite-graph statement is formalized as
\record{R002.finite_literal_path_decision}. Its hypotheses and fixed
start must be supplied by geometry. The noncomputable construction
is not a certified executable Python or C classifier. For fullerene
duals, the nonfacial-three-cycle exclusion in the received prose route
uses the established cyclic edge-connectivity result cited in
\cite{Connectivity}; the original proof of that external result was
not rechecked by the editor.

\subsection{Signatures and eventual validity}\label{def:signature}
\claimid{def:signature}
Work in the triangular lattice $\Lambda=\Z^2$, with
$R(x,y)=(-y,x+y)$. A signature specifies a fixed finite set of polygon
charts, paired boundary segments, corner equivalence classes, and
transition maps $z\mapsto R^jz+d$. Corner coordinates and translations
are integer affine functions of the parameters; rotations and gluing
combinatorics are fixed. Along a ray $p(t)=at+b$, $t\in\N$, we require
eventually positive edge lengths, embedded chart interiors, only the
specified seam identifications, noncrossing macro edges, and a valid
simple fullerene dual. Every cone has a marked corner label. Seams
need an ownership convention, and full corner equivalence classes are
used. Positive parameters and positive area alone do not prove validity.

A component $a_it+b_i$ eventually lies in the positive integers exactly
when $a_i>0$ or $a_i=0,b_i\ge1$. In the first case an entry bound is
$\max(0,\lceil(1-b_i)/a_i\rceil)$. Replacing 1 by 2 gives the analogous
bound for a lower limit two. A negative slope has no infinite positive
tail. These arithmetic entry bounds do not certify the quotient.

\section{The main result: convex lattice-face decompositions}\label{sec:main}
\begin{alignedtheorem}[Convex lattice-face rays]\label{thm:aligned}
For all sufficiently large integers $t$, suppose a valid simple
fullerene dual has a decomposition into a fixed finite number of
embedded closed convex lattice-polygon disks. Every side has positive
length and one of the six lattice directions. Straight angles and
auxiliary degree-six corners are permitted. No distinct boundary
points of one disk are identified. Corner coordinates and transition
translations are integer affine in $t$, while rotations, side pairings
and corner identifications are fixed. All degree-five vertices occur
among the marked corners, with a fixed label set. Then the set of
$t$ for which a literal path exists is eventually periodic.
The conclusion also holds when these hypotheses hold separately on
finitely many residue classes modulo a fixed integer.
\end{alignedtheorem}
\claimid{thm:aligned}
Thus some $t_0\ge0,P\ge1$ satisfy the same decision at $t$ and $t+P$
for all $t\ge t_0$. The received theorem proves existence. No practical
threshold, least period or implemented geometric-to-MSO/DFA compiler
is supplied. The specification-level existence theorem is now
Lean-proved; the prose correspondence lemma below transfers it to
these geometric hypotheses. Section~\ref{sec:proof} retains the
independent prose route; its original proof, corrections and
coordinator review are fixed in the packet. The hypotheses describe the objects
to which it applies, not a generic validity-testing algorithm.


\paragraph{Prose correspondence lemma (not formalized in Lean).}
The geometric hypotheses of Theorem~\ref{thm:aligned} determine a fixed
\record{ClaudeChallenge2.Sig} $S$ and a tail on which
\record{Valid} $S\,t$ and \record{TrianglesInCharts} $S\,t$ hold;
the quotient \record{graph} $S\,t$ is the stated fullerene dual and
\record{LiteralPath} is Definition~\ref{def:literal}.

\emph{Prose justification.}
For a convex lattice polygon with the six permitted side directions,
its supporting lines give exactly three closed strips
$L_x\le x\le U_x$, $L_y\le y\le U_y$,
$L_s\le x+y\le U_s$. Their extrema are selected from finitely many
affine corner coordinates; after their orders stabilize, all six
bounds are integer affine. Retain the fixed charts, rotation exponents,
affine translations and marked corners as the fields of $S$. For each
paired side, use its prescribed transition as a gluing record.

\emph{No extra identifications from unrestricted records.}
Let $g$ unfold a paired segment $A\subset K$ onto $g(A)\subset K'$,
and let $\ell$ be its supporting line. The consistently oriented disks
$g(K)$ and $K'$ occupy opposite closed supporting halfplanes, so
$g(K)\cap K'\subset\ell$. Their overlap there is an interval $I$
containing $g(A)$. If $I$ continues beyond an endpoint of $g(A)$,
both supporting edges continue straight at that point. The two disks
each contribute angle $\pi$. A fullerene dual has total cone angle
$5\pi/3$ or $2\pi$, so the point is flat, the two disks exhaust its
neighbourhood, and their actual seam must continue across it. The
orientation-preserving lattice isometry along the continued seam is
the same $g$: it agrees on a nontrivial line segment. Repeat across
the finitely many side subdivisions until reaching the endpoint of
$I$. Thus every pair identified by this unrestricted record is truly
the same surface point, even if it lies beyond the named subdivision.
Conversely, the prescribed paired sides are included in the records.
A self-pairing of distinct points is excluded by the embedded-closed-disk
hypothesis; an identity record is harmless. Transitive closure therefore
gives exactly the prescribed seam and full corner equivalences, without
extra identifications. This supporting-line and cone-angle argument is
prose, not a theorem certified by the Lean kernel. Thus sites
modulo these records are precisely the surface vertices, and common-chart
unit steps are precisely the graph edges. Closed-disk embeddedness
supplies the injectivity clause of \record{Valid}; the fullerene degrees
and marked-corner hypothesis supply its degree clauses.

Use the same external fullerene fact as in Section~\ref{sec:definitions}:
all dual three-cycles are facial, by cyclic edge-connectivity
\cite{Connectivity}. Each edge therefore has exactly the two distinct
opposite vertices of its incident triangular faces as common neighbours.
This gives \record{TwoCommonNeighbors}. Since chart boundaries follow
lattice edges, no open unit triangle is cut by a boundary; the covering
decomposition puts every triangular face inside a closed chart.
Every graph triangle is facial, so its three vertices lift there with
all three differences unit vectors. This is exactly
\record{TrianglesInCharts}. These geometric identifications and the
external fullerene fact are part of this prose bridge, not additional
Lean-proved declarations or a fresh audit of the cited proof.

Apply \record{ClaudeChallenge2.main} on the resulting tail. For fixed
residue branches write $t=r+qn$ and apply the construction to each
branch. The received Lean declaration \record{ClaudeChallenge2.main_residue}
then combines the branches in the original parameter, assuming eventual
agreement of the decision with each branch's \record{PathExists}.
That agreement is supplied by this prose representation lemma;
the combination of eventual periods is Lean-proved.
The owner confirmed the version-2 specification's meaning on
20 October 2026; the check is recorded in \record{OWNER-STATEMENT-CHECK.md}.
It is not proof peer review or review of this prose correspondence lemma.

\begin{catalogcor}[Specified catalogue families]\label{cor:catalog}
On every integer affine ray eventually remaining in the valid positive
integer domain of each of A1, A2, B2, C1, K3, E10, D12, D14, D1 and O2,
literal-path existence is eventually periodic. The same conclusion
holds for A3 on rays eventually satisfying $k>l>0$, by its separate
classification and arithmetic argument~\cite{A3}.
\end{catalogcor}
\claimid{cor:catalog}
There are eleven families in this statement. Positive boundaries with
a parameter fixed at 1 are included, notably K3 and O2 and E10 with
$r=1$. Collapsed zero boundaries are not. The corollary gives no
two-variable closed classification merely from ray periodicity.
The K3 part now also has a direct Lean proof on its quotient graph,
including the positive boundaries, without invoking the general prose
correspondence lemma; the other family routes retain their stated scope.
A3 generally has changing parent-edge directions and is not silently
placed under Theorem~\ref{thm:aligned}. The catalogue terminology and
corrections are Graver's \cite{Graver}; the constructions are not claimed
as new. Additional leapfrog families are discussed only as prospects
in Section~\ref{sec:extension}.

\section{Proof of the convex-face theorem}\label{sec:proof}
This section presents R004's second derivation on the common
six-direction domain. It does not use the uncompleted oblique-direction
transfer conditions. Its arithmetic-to-word reduction is specific;
arbitrary Presburger relations are not claimed to be finite-word MSO.

\subsection{Maximal face passages and their outputs}\label{lem:rows}
\claimid{lem:rows}
Rotate a passage direction to $(1,0)$. A suitable phase of the forced
vertex sequence is
\begin{equation}\label{eq:rows}
 q_{2j}=(j,h),\qquad q_{2j+1}=(j,h+1).
\end{equation}
A convex six-direction polygon is specified by
\[
 L_x\le x\le U_x,\qquad L_y\le y\le U_y,\qquad
 L_s\le x+y\le U_s.
\]
Its bounds are affine in $t$ after the finitely many max/min orders
stabilize. Redundant collinear sides do not alter this description.
With $L_y\le h\le U_y-1$, the triangle
$T_i=\{q_i,q_{i+1},q_{i+2}\}$ lies in the polygon precisely in the ranges
\begin{equation}\label{eq:triangle-bounds}
\begin{aligned}
 A_e&=\max(L_x,L_s-h),& B_e&=\min(U_x-1,U_s-h-1),\\
 A_o&=\max(L_x,L_s-h-1),& B_o&=\min(U_x-1,U_s-h-2).
\end{aligned}
\end{equation}
Even indices are $2j$ for $A_e\le j\le B_e$, odd indices are $2j+1$
for $A_o\le j\le B_o$. Put $\delta=A_e-A_o$ and
$\epsilon=B_o-B_e+1$, both in $\{0,1\}$. The full triangle-index set
is empty or the interval
\begin{equation}\label{eq:passage}
 \ell=2A_e-\delta\le i\le u=2B_e+\epsilon.
\end{equation}
This includes the one-odd-triangle case $A_e=B_e+1$,
$\delta=\epsilon=1$. It must not be discarded as an empty even range.

The first two vertices of an entered passage already belong to the
previous passage. Its new output is only
$q_{\ell+2},\ldots,q_{u+2}$, comprising
\begin{equation}\label{eq:output}
\begin{array}{ll}
 y=h:   &A_e+1\le x\le B_e+1,\\
 y=h+1: &A_e+1-\delta\le x\le B_e+\epsilon.
\end{array}
\end{equation}
For example the upper row has
$x\ge L_x+1-\delta$, $x+y\ge L_s+2-\delta$,
$x\le U_x-1+\epsilon$, and $x+y\le U_s+\epsilon$.
The last constant was corrected by hand in the received proof, followed
by a new finite check; the check did not originally discover that error.
Empty rows are allowed, but a nonempty passage has nonempty output.

Root and terminal passages can be truncated. There are finitely many
marked cone corners and at most ten oriented initial triangles at each.
The initial triangle fixes $h$ and an index of the form twice an affine
expression plus fixed parity. Store its first two vertices separately.
An endpoint corner and choice of output row similarly determine the
last index. Truncation adds affine bounds on $x$. Both cuts are made
when the whole path lies in one passage. Internal degree-five corners
are not compulsory stopping points. The two endpoint neighbours must
be distinct from all degree-five labels.

\subsection{States, transitions and all collisions}\label{lem:restricted}
\claimid{lem:restricted}
A state is $(\tau,h)$, where the finite type $\tau$ records face,
direction, max/min branches, parity, subdivided entry/exit side, and
root/terminal information. A constant $C$ bounds $|h|\le Ct+C$.
The number of states or face passages need not be bounded independently
of $t$. The relevant relations are Boolean combinations of
\begin{equation}\label{eq:restricted}
 \pm h\le at+b,\qquad \pm h\pm h'\le at+b,
 \qquad at+b\ge0,
\end{equation}
and fixed-modulus congruences. A relation such as $2h+h'\le at+b$
with distinct free variables is not part of this language. Substituting
one variable for another can be reduced by fixed parity and rounding.

Within a fixed type, each endpoint's three normal coordinates has the
form $c(t)$, $h+c(t)$ or $-h+c(t)$. When the last ordered pair of a
passage is glued to the next one's first ordered pair, the directions
are rotated together. In direction-aligned coordinates the linear part
of this transition is the identity, so $h'=h+c(t)$, not a doubling map.
The paired unit edge determines the opposite face; its interior
contains no macro corner. The conditions therefore have
form~\eqref{eq:restricted}.

Each output row fixes one of $x,y,x+y$ at $\pm h+c(t)$ and imposes
affine bounds in all three normals. For distinct row normals the
determinant is $\pm1$: their intersection is integral, and the third
normal is their sum or difference. Substitution into both rows gives
\eqref{eq:restricted}. Equal normals first require equal values;
the remaining integer interval is nonempty exactly when every lower
bound is at most every upper bound. The same restricted form results.
Between different faces, compare through every shared side and every
common marked corner label. Rotations permute the three normals with
signs. Embeddedness and fixed gluing ensure these exhaust quotient
identifications, even when faces share several sides. Auxiliary flat
corners are included. Comparison with the root's two saved vertices
uses the same formulas. Since an embedded face row is injective,
pairwise disjoint outputs and their disjointness from the root are
exactly the full vertex-distinctness condition.

\subsection{Folding the restricted arithmetic into a finite word}\label{lem:folding}
\claimid{lem:folding}
For $t\ge2$, write a height variable uniquely as $h=jt+y$, $0\le y<t$.
The quotient $j$ ranges over a fixed finite set. Set $t=2n+\eta$,
$\eta\in\{0,1\}$, and split $y$ into
\begin{equation}\label{eq:fold}
 L:y=v,\qquad R:y=t-1-v\quad(0\le v<n),\qquad
 M:y=n\quad(\eta=1).
\end{equation}
Original type, $j$ and $L/R/M$ are finitely many copies of the positions
of a unary word of length $n$. A middle copy is placed at the first
position and interpreted as coordinate $n$, independent of that position.

For two ordinary positions the inequalities become
\[
 \alpha v+\beta w\le 2cn+d,\qquad \alpha,\beta\in\{-1,1\}.
\]
The coefficient of $n$ is even. If $\alpha=-\beta$, then for $c\ne0$
truth stabilizes, while $c=0$ is order with a fixed successor offset.
If $\alpha=\beta=1$, the sum lies between 0 and $2n-2$. The only
nontrivial cases are $c=0$, a finite list near the left end, and $c=1$.
For $c=1$, putting $a=n-1-v,b=n-1-w$ gives
$a+b\ge-d-2$, either automatic or the exclusion of finitely many
near-end pairs. The case $\alpha=\beta=-1$ is its integer complement.
If one position is a middle copy, only one variable remains and the
nontrivial inequalities are fixed distances from an end. Two middle
copies give an affine truth value that stabilizes.

Fixed-modulus congruences are represented by cyclic MSO colours on
the word: the first position has colour zero and each successor
increments the colour. The final colour determines $n$ modulo the
fixed modulus. Finite copies of element or set quantifiers are expanded
copy by copy. Thus every fixed Boolean combination used above is
MSO-definable. This argument does not permit an arbitrary Presburger
relation such as $v+w\le n$; it uses the even coefficient obtained by
the folding step.

\subsection{The exact path sentence and unary regularity}\label{lem:mso-path}
\claimid{lem:mso-path}
Choose a nonempty state set $S$, root $a$ and terminal $b$. Terminal
states have no outgoing transition and root states no incoming one.
Record the parity of the number of vertices already output as a finite
type. Initially it is even, since the root has two saved vertices;
each transition adds the parity of $u-\ell+1$.
Impose the following MSO conditions:
\begin{enumerate}
\item If $a=b$, $S=\{a\}$. Otherwise $a$ has in/out degrees $(0,1)$,
      $b$ has $(1,0)$, and every other state has $(1,1)$ within $S$.
\item Every nonempty proper subset of $S$ has a transition, in either
      direction, across its boundary. This is unary-set connectivity.
\item Outputs of distinct states are disjoint, and all outputs are
      disjoint from the two saved root vertices.
\item Endpoint degrees are $5,6,6,5$ as in Definition~\ref{def:literal}.
      Total vertex count is even and at least four.
\end{enumerate}
For one state, the last size condition is $u-\ell+1\ge2$; with two
or more nonempty-output states it is automatic. Degree conditions are
membership or nonmembership in the finite marked-cone set. No binary
relation quantifier is used.

On a finite set, the first two conditions give one directed path
through all of $S$: a reached cycle would create an extra incoming edge,
and a separate cycle would contradict connectivity. Join the passages
in this order. Ordered seam pairs give consecutive and skip adjacency;
the third condition gives distinct vertices. Size, parity and degree
conditions give exactly the literal path.

Conversely, split any literal path into maximal consecutive passages
in its convex faces, truncating only at its ends. Equations
\eqref{eq:triangle-bounds}--\eqref{eq:output} recover their states.
Nonempty output and path simplicity prevent a repeated state. Any extra
eligible successor within $S$ would share the next output vertex with
the actual successor, contradicting disjointness. Hence the selected
states satisfy all four conditions. One-passage paths and paths with
internal cones are included in both directions.

Finite-word MSO defines exactly the regular languages
\cite[Section 3.1, Theorem 3.2]{MSO}. Apply this established result
only at this last logical step. For each $\eta$, a deterministic
automaton over one letter eventually repeats a state, so its accepted
lengths $n$ are eventually periodic. Combine $t=2n+\eta$ and, if
necessary, the initially fixed residue branches. This proves
Theorem~\ref{thm:aligned}. The geometric reduction belongs to the
received R004 proof, not to the cited automata theorem. Neither the
full proof of that external theorem nor an MSO compiler was newly
verified by this manuscript's editor.

\section{Applying the theorem to the specified families}\label{sec:families}
\label{prop:family-decompositions}\claimid{prop:family-decompositions}
The applicability evidence is R004's fixed family-decomposition report
and its coordinator review. D12 uses two subdivided hexagons and D14
two convex hexagons whose long sides are subdivided with lengths
$r,s,r$. O2 uses two hexagons and six triangles. K3 uses two caps,
each divided into four triangles, and three hexagons. These
decompositions retain the positive parameter-1 boundaries.

D1's original thirty-degree sides do not directly satisfy the six
direction hypothesis. Cut a triangle $A_iB_iW_i$ from each oblique side
and move it to the back by the same lattice isometry. For example,
\[
 A_0=(r,0),\quad B_0=(r+p,p),\quad W_0=(r,p),\quad C=(-p,r+2p),
\]
and $W_0-C=\alpha(A_0-C)+\alpha(B_0-C)$ with
$\alpha=(r+p)/(2r+3p)$ and $0<2\alpha<1$. The pieces are contained
and disjoint. The result is eight convex pieces: one regular hexagon,
six quadrilaterals, and the back hexagon. New corners are flat degree
six. Area equality alone is not used to infer the same quotient.

For E10 put $m=2r+s$, tube circumference $C=3m$. Move the longitudinal
cut from $(p,p)$ to $(0,2p)$, obtaining three lattice parallelograms
of width $m<C$ and two caps: five disks, also when $r=1$. Cut-corner
positions involve $p\bmod C$. If $C$ grows, $\lfloor p/C\rfloor$
eventually stabilizes; if $C$ is constant, use fixed residue classes.
Delete duplicate auxiliary corners and fix their order to obtain
affine coordinates and fixed gluing. This does not classify E10 with
$r=1$.

\begin{center}\small
\begin{tabular}{p{.19\textwidth}p{.69\textwidth}}\toprule
Family and domain & Received reason for all valid affine rays\\\midrule
A1, $r\ge1$ & Integer GC absence, including the separate $r=1$ case; eventual period one.\\
A2, $p\ge1$ & Nonconstant rays eventually give nonprimitive GC images; constant rays are constant graphs.\\
B2, $p,r\ge1$ & Its closed condition is a finite disjunction of affine equalities.\\
C1, $a,b\ge1$ & Its closed condition is $b=1$.\\
A3, $k>l>0$ & Separate complete classification and arithmetic; Section~\ref{sec:arithmetic}.\\
K3, $r,s\ge1$ & Lean literal-path ray periodicity, including parameter 1; no general prose bridge is used.\\
E10, $p,r,s\ge1$ & Arithmetic for $r\ge2$; the five-piece decomposition includes $r=1$.\\
D12, D14, D1 & The convex decompositions on valid positive integer parameters.\\
O2, $t,d\ge1$ & Convex decomposition, including parameter 1; absence is received only for $t,d\ge2$.\\\bottomrule
\end{tabular}
\end{center}
This proves Corollary~\ref{cor:catalog} using the stated geometric
receipts and the separate arithmetic results. It does not cover all
112 catalogue entries or every zero-boundary degeneration.

\section{Closed conditions and the scope of partial results}\label{sec:conditions}
The following are received classifications of valid catalogue graphs.
The R004 prose arguments are retained; K3 now additionally has a full
Lean graph classification on its specified closed domain. Independent derivation of the same conclusion is distinct
from checking every coordinate convention in the research team's
tables. Lean arithmetic for a predicate does not by itself prove its
equivalence to a graph property.

\subsection{The specified closed domains}\label{prop:closed}
\claimid{prop:closed}
\begin{center}\small
\begin{tabular}{p{.19\textwidth}p{.36\textwidth}p{.31\textwidth}}\toprule
Family/domain & Literal path condition & Received scope\\\midrule
A1, $r\ge1$ & Never positive & Integer GC for $r\ge2$; at 1 no degree-six vertex. R004 report 4 and GC proof.\\
A2, $p\ge1$ & $p=1$ & Existing $C_{60}$ witness and integer GC. R004 report 14 does not supply a new $p=1$ witness.\\
C1, $a,b\ge1$ & $b=1$; IPR iff $a,b\ge2$ & R004 report 3 and C1 proof; coordinator review of the team's argument.\\
B2, $p,r\ge1$ & $p=1$ or $|p-r|=1$ & R004 report 4 and B2 proof, whole positive domain.\\
K3, $r,s\ge2$ & $s\equiv\pm1\pmod r$ & R004 positive and necessary prose proofs retained; full Lean graph iff now received, without a prose bridge.\\
E10, $p,s\ge1,r\ge2$ & $m=2r+s$ and $p\equiv d\pm1\pmod m$ for some $d\in\{0,r,r+s\}$ & R004 reports 5--6 and both prose directions. Predicate-ray arithmetic formalized; full graph iff not formalized.\\\bottomrule
\end{tabular}
\end{center}
K3 writes $s=qr+u$: $u=0$ is excluded by integer GC,
$2\le u\le r-2$ by belts of widths $r,u,r-u$, and $u=1,r-1$
by explicit simple paths. The overlapping small cases $r=2,3$ are
separate. The iff is not extended to $r=1$ or $s=1$. E10's candidate
congruences come from the positions of corners around the tube;
reaching a cone does not alone establish simplicity. The research
team's full table of eighteen E10 regions, walls and witnesses is not
all independently reviewed.

\subsection{Partial branches and boundary examples}\label{prop:partial}
\claimid{prop:partial}
O2 is received as absent and IPR for $t,d\ge2$, with the mapping-metadata
disclosure retained. For D14 and D1 with both parameters at least two,
the research team has belt-based absence drafts; R004's all-ray
periodicity receipt is not a new independent proof of that closed
absence claim. Common integer scaling by at least two gives absence
in the specified leapfrog domains, but does not classify their
primitive points.

For E10$(p,1,s)$, $m=s+2$, the research team's witnesses give
sufficiency when $p<m$ or
$p\bmod m\in\{0,1,2,m-2,m-1\}$. General necessity is not proved.
Among the 84 positive $(p,s)$ with
$N=12+2s^2+12s+12p(2+s)\le740$, exactly E10$(9,1,4)=C_{740}$
is negative; all 83 smaller examples are positive. E10$(17,1,11)$
and $(31,1,7)$ give further negatives at $C_{3038},C_{3542}$.
These three examples are non-IPR. R004's report 4 independently
constructed and checked them. This is minimality only inside the
$r=1$ E10 family, not a general iff.

The pre-existing D12 branches are retained as historical scope:
$\gcd(r,s)\ge2$ gives absence; $r\ge3$, $s=qr\pm1$, $q\ge1$
gives existence by the separately derived simple witnesses. The
research team's $|r-3s|=1$ and $r\ge4s$ branches have their earlier
review limitations. A full D12 classification, including newer work,
is outside this paper and this packet.

\section{Arithmetic on affine rays and exact examples}\label{sec:arithmetic}
\subsection{The determinant reduction}\label{lem:divisibility}
\claimid{lem:divisibility}
Let $A(t)=\alpha t+\beta$ be eventually positive and
$B(t)=\gamma t+\delta$. A constant $A$ gives a fixed-modulus
congruence. If $\alpha>0$, then
\begin{equation}\label{eq:determinant}
 A\mid B\quad\Longrightarrow\quad A\mid\alpha\delta-\gamma\beta.
\end{equation}
For nonzero right side, divisibility is impossible once
$A>|\alpha\delta-\gamma\beta|$. For zero right side, the ratio
$B/A=\gamma/\alpha$ is constant and is either always integral or
never integral. Finite Boolean combinations of affine divisibilities
are eventually periodic; a growing modulus gives eventual constancy.
Take a common multiple of fixed moduli and a maximum stabilization
threshold. This elementary arithmetic is not proposed as new number
theory. The received Lean files establish the exact predicates below.

For all integer ray coefficients, define parameters by
$(at+b).\operatorname{toNat}$. The formal statements for all rays use
\emph{domain and positive predicate}; they do not call out-of-domain
graphs positive. On an eventually positive domain the conversion
agrees with the integer affine parameter.
\label{thm:k3-arithmetic}\claimid{thm:k3-arithmetic}
\record{Claude2.k3_ray_all} proves eventual periodicity of
$\mathrm{K3Dom}\wedge\mathrm{K3Pos}$ on all integer affine rays,
where the domain is $r,s\ge2$ and the predicate is
$r\mid(s-1)$ or $r\mid(s+1)$. If $r$ is nonconstant on a valid tail,
\record{k3_ray_const} gives eventual constancy. If $r=b\ge2$ is
constant, \record{k3_ray_period} gives period $b$ from any specified
domain-entry threshold. These need not be least periods.
\label{thm:e10-arithmetic}\claimid{thm:e10-arithmetic}
\record{Claude2.e10_ray_all} proves the corresponding all-ray statement
for domain $p,s\ge1,r\ge2$ and the six divisibilities in the E10
condition. If $2r+s$ is nonconstant, the predicate is eventually
constant; when it is constant, \record{e10_ray_period} gives that
positive modulus as a period on the given domain tail.
\label{thm:other-arithmetic}\claimid{thm:other-arithmetic}
\record{b2_ray_all}, \record{c1_ray_all}, \record{a1_ray_all} and
\record{a2_ray_all} give eventual constancy of their domain-qualified
predicates. The conjunction/disjunction/negation closure and domain
qualification are proved in the same arithmetic layer.

\subsection{A3 is a separate variable-direction case}\label{thm:a3-arithmetic}
\claimid{thm:a3-arithmetic}
For $k(t)=at+b,l(t)=ct+d$ set $\Delta=ad-bc$. If $\Delta\ne0$,
$\gcd(k,l)=\gcd(k,l,\Delta)$, so primitivity is periodic modulo
$|\Delta|$. When $\Delta=0$ and the ray is nonconstant, it is
eventually an increasing integer multiple of a fixed primitive
direction and hence negative by integer GC. Constant rays are constant.
In the primitive case put $m=\min(l,k-l)$. The A3 criterion is
$m=1$, or quotient $q=\lfloor k/m\rfloor\ge3$ with $m\mid(k\pm1)$,
or $q=2$ with $r=k-2m=1$ or $r\mid(m\pm1)$. Branches and quotient
stabilize, including an integer limiting ratio by checking the sign
of the constant remainder. The persistent $q=2,r=0$ growing branch
is not primitive. Apply~\eqref{eq:determinant} to the remaining tests.

The first paper proves the full graph classification. Its
\record{ClaudeCor.criterion_ray}, restated as \record{Claude2.a3_ray},
formalizes the arithmetic existence of a threshold and positive
period. A separate release bridge gives graph-path periodicity.
Neither result certifies an executable threshold/period algorithm.
R004's separate arithmetic argument permits period $|\Delta|$ for
nonconstant $\Delta\ne0$ rays and period 1 otherwise. This is not a
least-period statement or formalization of the general geometric
theorem. After classification, eventual periodicity is also a
consequence of \cite[Theorem 1.15]{BGW} and is not a separate novelty claim.

\subsection{Four previously solved rays}\label{ex:exact-rays}
\claimid{ex:exact-rays}
\begin{center}
\begin{tabular}{lll}\toprule
Ray, $t\ge0$ & Positive values & Valid period and threshold\\\midrule
D12$(124+3t,41+t)$ & All $t$ & $P=1,t_0=0$\\
D12$(125+3t,41+t)$ & None & $P=1,t_0=0$\\
K3$(43+2t,170+5t)$ & $40,42$ & $P=1,t_0=43$\\
E10$(200+5t,17+2t,23+t)$ & $3,10,17,34$ & $P=1,t_0=35$\\\bottomrule
\end{tabular}
\end{center}
The original compressed-word calculations did not use the known
family congruences. Later independent prose arguments recover all
$t$ from independently derived family conditions, or from explicit
D12 blocks. The expected statements were known, not blind predictions.
\label{ex:k3-formal}\claimid{ex:k3-formal}
\record{Claude2.k3_example} now proves for every $t\in\N$
$\mathrm{K3Pos}(43+2t,170+5t)\leftrightarrow(t=40\lor t=42)$.
\label{ex:e10-formal}\claimid{ex:e10-formal}
\record{Claude2.e10_example} proves the corresponding E10 predicate
iff $t\in\{3,10,17,34\}$. The determinant constants used are
$123,127$ for K3 and $142,144,739,749,681,691$ for E10. Divisibility
bounds $t$ before the remaining cases are checked by kernel arithmetic.
These are all-$t$ arithmetic proofs, not graph iff proofs by themselves.

The earlier D12 two-ray receipt covers every $s\ge41$:
D12$(3s+2,s)$ is negative and D12$(3s+1,s)$ has a path with
$56s+12$ distinct vertices. Its reported computer-assisted proof checks
2,058 nonnegativity conditions and 756 no-integer-solution conditions,
2,814 in total. A separately coded reread by the same verifier is not
a further reviewer. The coordinator's independent construction at
$s=58$ agrees, with a 3,260-vertex positive witness; it does not audit
all 2,814 conditions. These records are neither Lean proofs nor a
full D12 classification.

\section{Received formalization and what it does not cover}\label{sec:formal}
\subsection{Theorem 1a on the coordinator's specification}
The current entry points are \record{ClaudeChallenge2.main} in
\record{Formal/Claude3ComparatorSolution.lean} and
\record{ClaudeChallenge2.State.main_statement} of type
\record{ClaudeChallenge2.MainStatement} in \record{Formal/Claude3Thm1a.lean}.
For every fixed signature $S$ and $t_1$, they prove that if
\record{Valid} $S\,t$ and \record{TrianglesInCharts} $S\,t$ hold for all
$t\ge t_1$, then some $t_0$ and positive $P$ satisfy
\record{PathExists} $S\,t\leftrightarrow$ \record{PathExists} $S\,(t+P)$
for every $t\ge t_0$. The independent challenge was written by the
coordinator before the proof and amended to version 2 to include
\record{TrianglesInCharts}. The paper-to-specification passage is the
prose lemma following Theorem~\ref{thm:aligned}; it is not formalized.
The Lean assumptions are weaker than the paper's geometric hypotheses:
they do not require a sphere, exactly twelve cones, or nondegenerate
polygon disks. With the same conclusion, the exact formal result is
therefore a stronger sufficient theorem, not a literal restatement
of the paper's geometric objects.
\record{ClaudeChallenge2.main_residue} in \record{Formal/Claude3Residue.lean}
proves the fixed-residue variant after the decision-to-branch agreements
are supplied.

The proof was supplied by Claude workers, through finite-word MSO,
height folding, chart geometry, quotient identifications, state
reduction and sentence assembly. It is not an independent Lean
proof by the research team. The research team's serial build and
fresh axiom queries for all three declarations found only
\record{propext}, \record{Classical.choice} and \record{Quot.sound}.
The saved coordinator Comparator receipt reports six kernels
accepting the proof of \record{ClaudeChallenge2.main} in paranoid mode;
that receipt is not a separate six-kernel check of the residue declaration.
This edit does not rerun external kernels or provide OS process isolation.
The owner's version-2 statement-meaning confirmation is recorded
separately. Human proof peer review remains unrecorded. Neither this formalization nor
R004's separate prose proof establishes novelty or an executable
geometric compiler, practical threshold, or least period.

\subsection{K3 graph classification and ray periodicity}\label{thm:k3-sufficient}
\claimid{thm:k3-sufficient}
The new entry declarations are
\begin{quote}\small\noindent
\record{Claude4K3.k3_statement} : \record{Claude2Aligned.K3Statement};\\
\record{Claude4K3.noThree};\\
\record{ClaudeChallengeK3.main}.
\end{quote}
For natural $r,s\ge2$, the first proves the full iff between an actual
literal path in \record{Claude2Aligned.k3Graph}$(r,s)$ and
$s\equiv\pm1\pmod r$. The second proves that no edge has three common
neighbours, using a classification of graph triangles in the cap or tube.
Together with the earlier necessity machinery, this supplies the missing
necessity direction. R004's positive and necessary prose proofs remain
separate evidence.

The third proves eventual periodicity of literal-path existence on every
ray $(r,s)=(r_0+at,s_0+bt)$ with $r_0,s_0\ge1$ and $a,b\ge0$.
There are no \record{Valid} or \record{TrianglesInCharts} hypotheses.
The K3 gap previously filled by the general prose correspondence lemma is
therefore closed in Lean on this explicit quotient specification. The
proof constructs the needed common-neighbour property directly and does
not construct a signature isomorphism into Theorem 1a.
Nonconstant parameter-one boundary rays use explicit eventual paths;
constant rays are periodic because their decision never changes. The
modular iff itself is still restricted to $r,s\ge2$.

The specification uses eleven convex polygons, equivalence closure of
explicit affine side gluings, and adjacency from unit differences inside
a common chart. The literal path and actual neighbour-set degree
 definitions are those of the first paper's specification.
The earlier sufficient declaration \record{Claude2K3.k3_sufficient}
of type \record{Claude2Aligned.K3Sufficient} is retained.

The proof follows R004's tube construction. A glued-site invariant
is $(\mathrm{hgt},\mathrm{ang}\bmod6r)$. It is preserved by all
42 directed edge gluings. It is not globally injective: the two
central cap triangles are not separated by it. On tube points with
horizontal coordinate difference at most one, however, equality of
quotient vertices forces equality of points. The explicit paths are
\[
 P_n^+=(-(n\bmod2),\lfloor(n+1)/2\rfloor),\qquad
 P_n^-= (n\bmod2,\lfloor n/2\rfloor),
\]
with $4s-2$ vertices for $s=qr+1$ and $4s+4$ for $s=qr-1$.
Chart membership, consecutive and skip adjacency, distinctness, and
the four endpoint degree requirements are proved on the quotient.
The invariant is not used to claim a global vertex normal form.

\label{lem:k3-corners}\claimid{lem:k3-corners}
\record{Claude2K3.deg_corner} proves degree five for all twelve labelled
corner representatives. It does not by itself prove they exhaust all
degree-five vertices or all global vertex types.
\label{lem:k3-hex-interior}\claimid{lem:k3-hex-interior}
\record{deg_hex_interior} proves degree six in hexagon interiors.
\label{lem:k3-tri-interior}\claimid{lem:k3-tri-interior}
\record{deg_tri_interior} proves degree six in the six noncentral cap-triangle interiors.
General seam interiors and the central cap triangles are not covered
by those older partial degree lemmas alone. The new necessity proof
uses the cap/tube triangle classification and common-neighbour bound;
the K3 iff in Section~\ref{sec:conditions} is now formalized. This does
not relabel the partial degree lemmas as a global degree classification.

The older graph definitions were written by a Claude worker. Before
the new proof, the coordinator copied them into a separate K3 challenge.
Static comparison confirms the copied definitions and the version-2
challenge/solution statement agree. An independent AI reader of version 1
found the missing parameter-one boundaries; version 2 corrected them
before the proof. The reader's finite checks and a separate Python mirror
support the catalogue correspondence, without constituting a universal
catalogue-isomorphism proof or human proof review. No owner-confirmation
field or procedure is used for K3. Shared specification files contain
unused D12 definitions, but no new D12 classification is a root or
claimed result of this package.

\subsection{Historical conditional kernels}\label{sec:historical-formal}
\label{lem:historical-kernels}\claimid{lem:historical-kernels}
The earlier research layer supplies conditional pieces:
\record{finite_literal_path_decision} assumes a finite type and two
common neighbours; GC lattice transport excludes odd endpoints in a
scaled lattice; the four cylinder arithmetic declarations prove the
height-kernel displacement with stated divisibility hypotheses;
\record{induction_uniform_depth} assumes a finitely branching child
list and genuinely decreasing lexicographic rank. The concrete
\record{ietComponentChildren} models only removal of equal last labels,
not every subtraction, equal-length elimination or accelerated suffix.
The geometric hypotheses, all-start coverage, full compressed action
and a certified program are not supplied by those conditional types.
There were historically 57 such research declarations; the former
total 98 also counted 41 A3 arithmetic entries. Neither count is a
formalization percentage or the current library size.

\subsection{Pins, receipt scope and proof dependencies}\label{sec:formal-records}
The received toolchain is Lean 4.35.0 with the fixed mathlib manifest.
The earlier arithmetic and K3 sufficiency reports record only
\record{propext}, \record{Classical.choice}, \record{Quot.sound} as
transitive axioms. Their receipt bytes and scopes are retained.
For Theorem 1a, its separate source closure and the earlier research-team
build/axiom output remain fixed. For the new K3 results, the coordinator
records a successful build and standard-three-axiom queries for all three
entry declarations. The research team receives and hash-checks these
records without running Lean.

The saved six-kernel K3 Comparator receipt checks
\record{ClaudeChallengeK3.main} in paranoid mode only; it is not a separate
six-kernel check of \record{k3_statement} or \record{noThree}.
It is also distinct from the earlier Theorem 1a and A3 receipts.
All six kernels check the same proof, not six independent derivations.
The K3 execution lacked OS process isolation on macOS. No external kernel
or Lean command is rerun for this edition.
The claims' expository dependencies and Lean's imported proof-term
dependencies are different. A prose statement can remain open for
its coordinate or human-semantic bridge while a separate exact
formal declaration is proved. Conversely the proved arithmetic
predicates do not supply their graph iff interpretation. No human
peer review of the proofs is recorded and novelty is not established.

\section{An extension to fixed oblique directions (proof sketch)}\label{sec:extension}
\label{claim:oblique}\claimid{claim:oblique}
\textbf{Status.} The coordinator has reviewed the skeleton of AN
Sections 2--6 as a plausible proof sketch. Explicit Section 5 transfer
conditions and numerical validation remain incomplete; no independent
second proof covers this wider domain. Four obligations below remain.
This section is not part of the main theorem or the abstract's result.

The proposed extension replaces the six side directions by an arbitrary
fixed finite set of primitive lattice directions, keeping eventual
validity, ownership and fixed affine gluing. For a direction $v$, use
transverse coordinate $h=\det(v,z)$. Across a seam the shift is affine.
Let $D$ be a common multiple of the nonzero $|\det(v,e)|$ and put
$Q=4D$. Add the transverse cell modulo $Q$ and $t$ modulo $Q$ to the
finite labels. Entrance and exit indices are
$k_{\rm in}=\lceil\xi_{\rm in}\rceil$,
$k_{\rm out}=\lceil\xi_{\rm out}\rceil$ with
$k_{\rm in}\le k<k_{\rm out}$. Reordering cells $n=\sigma+Qq$
conjugates the discrete midpoint permutation; it does not claim an
order-preserving conjugacy of the continuous intervals.

The AL/AM/AH proposal keeps a full partition into return towers,
including saved periodic components. Equal final labels remove a
periodic tower; equal lengths merge and remove a label without losing
an orbit. Repeated winner cycles are compressed as powers. The claimed
bound is on the number of Letter/Cat/Power nodes, not on expanded
passage count. Base widths and exponents are affine on fixed residue
branches; word lengths and displacement vectors are rational
polynomials in $t$. A lexicographic measure uses leading coefficient
sum in the original parameter, number of labels, and constant length;
it is not silently reused after coefficient-amplifying reparameterization.

For local ownership, close the side directions under rotations. A
unit triangle meets at most
$\sum_e(\max_\Delta\det(e,z)-\min_\Delta\det(e,z)+1)$
candidate macro-edge levels. Embedded noncrossing edges and finitely
many corner displays bound the length of the local gluing word.
A finite number of charts alone does not bound how often they can
be revisited. Entrance-closed, exit-open ownership excludes repeated
representations of one position from genuine quotient collisions.
Keep the full corner class and the degree-six second-point test;
non-IPR graphs can have fewer than 120 legal starts.

Tower return restores midpoint and direction; two returns restore
left/right phase. Nonempty output, initial suffix and terminal prefix
must be connected before using this as a finite first-repeat bound.
Positions are leaf types and ancestor-power repetition indices.
A free index $j$ contributes $j$ times a displacement depending only
on $t$. Products $jt$ or $jt^2$ are allowed coefficients in one-parameter
Presburger arithmetic; products of two freely quantified repetition
indices are not used. Fix one start address $\theta$ throughout:
\[
\begin{split}
\mathrm{Bad}(t,\theta,e)\ \Longleftrightarrow\
\exists\xi,\eta\,[D_\theta(\xi)\wedge D_\theta(\eta)
\wedge\xi\le e\wedge\eta\le e\\
\wedge\neg I_\theta(\xi,\eta)\wedge G_\theta(\xi,\eta)].
\end{split}
\]
Success would be $\exists\theta,e\,[\mathrm{Legal}(t,\theta,e)
\wedge\neg\mathrm{Bad}(t,\theta,e)]$. Legal includes odd endpoint index,
minimum length, degrees and the preceding output position. Allowing
each collision a different start address would be incorrect.

The remaining obligations are: (i) write the first-crossed-edge,
containment and order inequalities of the Section 5 transfer, with an
oblique-family example; (ii) verify the AH/AM affine tower coefficients
for the $Q$-split IET; (iii) keep effective threshold/period computation
dependent on the effective elimination result \cite{QE}, whose full
proof was not reaudited; (iv) compare an oblique-family ray's predicted
tail against full graph decisions. None is newly resolved here.
BGW supplies eventual periodicity only after the geometric formula
has been established.

\subsection{Prospective applications only}
C2 and K4 are specified leapfrogs of C1 and K3; the identified O1$(p,0,s)$
subfamily is a specified leapfrog of O2. The local coordinate map
$(x,y)\mapsto(x-y,x+2y)$ has determinant 3. These are prospects through
the extension, not families in Corollary~\ref{cor:catalog}. Leapfrog
is not assumed to preserve absence in general, and the full
three-parameter O1 family is not classified. These prospects do not
cover arbitrary signatures or every degenerate boundary. For the prospective signatures, twice area is $24a^2+48ab+12b^2$ for C2 and $42r^2+36rs$ for K4. Their $a=0$ and $r=0$ cases collapse faces, while $b=0$ and $s=0$ give zero sides. The zero middle coordinate in the identified O1 subfamily denotes a separately identified leapfrog signature, not permission to discard a zero edge in the old signature.

\section{Fixed-width structure, finite checks and novelty}\label{sec:evidence}
\subsection{Belts and a fixed cylinder}\label{prop:cylinder}
\claimid{prop:cylinder}
Developed triangle midpoints lie on a straight line; vertices alternate
on its two neighbouring integer-height rows. Under $z'=Mz+q$, a height
$h(z)=n\cdot z+e$ transforms as $n'=nM^{-1}$,
$e'=e-n'\cdot q$. If every start is covered by a return preserving
position, direction and primitive height, no cone lies inside a band,
and all widths are at least two, the required pair of endpoint rows
cannot both occur. Width one alone is not sufficient: simplicity and
endpoint-neighbour degrees remain necessary. For integer GC$(k,0)$,
$k\ge2$, the odd-point determinant $\pm1$ contradicts membership in
the $k$-multiple lattice; rotations about parent vertices preserve that
scaled lattice.

For the fixed-cylinder argument AD let circumference vector $W=gu$,
$g=\gcd(|W_1|,|W_2|)$, $h(z)=\det(u,z)$, and insertion vector $D$
have $H=h(D)>0$. Let $L$ be the least common multiple of the nonzero
$|h(v)|$ over the six unit directions. An entrance-state period is
\begin{equation}\label{eq:cylinder}
 Q=2gL,\qquad {Q\over h(v)}(Hv-h(v)D)\in\Z W.
\end{equation}
Directions with $h(v)=0$ loop in the flat core and repeat without
reaching a cone. If the caps and collars have a fixed finite vertex
set $C$, at most $|C|+1$ new cap visits occur before stopping. Overlap
of the same three-vertex display is not a new visit. On a residue
branch the passages have form $a_i(n)+jv_i$, $0\le j\le\ell_i(n)$,
with affine $a_i,\ell_i$ and fixed $v_i,W$. All collisions are
$a_i(n)+jv_i=a_k(n)+qv_k+zW$, with distinct-position and range
conditions. Ordinary Presburger arithmetic then gives periodicity.
$Q$ is an entrance period, not necessarily the final or least period.

The uniform cap argument does not extend to growing width/caps.
The old D12$(3,3t+4)$ example has passage word $(BA)^{t+1}$ and first
legal endpoint index $30t^2+88t+59$. Its $2(t+1)$ passages are unbounded.
The former general bounded-passage forecast was withdrawn.

\subsection{Dated finite evidence}\label{sec:finite}
\claimid{sec:finite}
The packet retains, rather than reruns, the relevant evidence.
T0770 compares 1,809 arithmetic values on nine known A3 rays;
T0771 adds fourteen graph decisions on an independent constructor.
Those fourteen are not unseen tests and the 1,809 values are not counted
twice. T0887 checks fixed tail coefficients beyond 64 on those rays.
T0888 audits E10's fixed descriptor: eight charts, 36 inverse gluings,
108 intervals, 156 connections, 120 descriptors and twelve labels.
T0889 audits 71 charts, 342 edges and 21 single-zero boundaries.
These are finite checks of supplied data, not the general ownership
theorem. The recorded AM audit has 6,564 inputs and 67,312 states;
the AN phase audit has 48 directions, 14,904 entrances, 4,005 intervals
(1,013 empty), 5,496 lifts and 40,380 conjugacy cells.

R004's report-14 finite arithmetic check has 378 coefficient
representations and 13,430 pairs, maximum recorded period 44 and
threshold 69. The report-15 checks include 7,056 strips, 2,255 nonempty
row outputs, 16,212 subintervals, 99,225 ordered row pairs and 509,160
folded-point comparisons. Their hand-found row correction and earlier
type/key failures are kept with their chronology. Numerical agreement
does not replace either general proof or repair its hypotheses.
No graph construction, R004 checker, or historical scientific audit
was rerun for this translation.

The K3 worker's earlier finite mirror agreed with its specification
and the verifier's construction on the stated known samples. The two
constructions and generated Lean text are separate from the kernel
proof. Generated degree formulas were fitted from sample data, but the
received Lean theorem checks their universal statements; samples alone
would not justify sufficiency. The worker's mirror is not a new
independent reviewer of its own specification.

\subsection{Novelty and remaining obligations}\label{sec:novelty}
\claimid{sec:novelty}
\paragraph{Immediate known consequences.}
Once a fixed sentence $\psi_\eta$ describes success on $a^n$,
$t=2n+\eta$, MSO/regular-language equivalence gives a finite automaton;
its unary state orbit is eventually a cycle \cite[Section 3.1]{MSO}.
The HTML numbering is Theorem 3.2; the same v1 PDF calls it Theorem 14.
Finite fixed-modulus residue branches preserve eventual periodicity.
Likewise, a fixed-size parametric Presburger description gives the
conclusion directly by \cite[Theorem 1.15]{BGW}. Neither theorem builds
the geometric description: an unbounded list of path vertices is not
a fixed list of quantified variables, and variable-variable products
are not permitted by that Presburger theorem.

For a fixed-width construction $G_n=C[F^n(B)]$ with fixed finite
boundary, literal-path existence has a fixed CMSO$_2$ description.
Quantify a vertex set $X$, an edge set $E'$ and distinct endpoints
$a,b$: $(X,E')$ is connected, has degree one at $a,b$ and degree two
elsewhere; $|X|\ge4$ is even; $a,b$ have ambient degree five and their
path neighbours ambient degree six; every two-edge path in $E'$ has
adjacent outer vertices in the ambient graph. This permits interior
cones and extra chords. The selected subgraph is precisely a literal
path, in both directions. Finite-state recognizability for a fixed
HR signature \cite[Chapter 1, Theorems 1.12 and 1.21(2)]{Courcelle}
then makes iteration of $F$ a finite-state orbit. Thus this fixed-width
existence statement is a short known consequence, not a new principle.
Mere bounded treewidth of an arbitrary indexed sequence is insufficient:
the uniform repeated construction is essential. The explicit $Q$,
cap-visit bound and actual period-calculation procedure require
separate evaluation.

\paragraph{Nearby results and the missing translations.}
Behmaram--Boutillier's fixed-circumference matching transfer matrix
has $2^m$ states \cite[Section 3, Eq.\ (3.1)]{Behmaram}; increasing
width does not leave the same matrix, and its observable is matchings,
not literal paths. Fischer--Makowsky prove periodicity or semilinearity
of model-size spectra for CMSOL classes with bounded width
\cite[Theorems 1.4, 1.7 and 1.10]{Fischer}. Such a spectrum asserts
existence of some model of a size, rather than deciding a specified
$G_t$. Kotek--Makowsky's labelled (bi-)iterative graph families and MSO
polynomials admit finite matrix transfers and $C^2$-finite recurrences
\cite[Definitions 17--20, Lemma 28, Theorems 30 and 32]{Kotek}.
Their bounded-clique-width framework does not automatically include
growing lattice patches; a recurrence for values is not itself
Boolean periodicity. In each case the geometric family and its
parameter require a uniform interpretation satisfying the cited
hypotheses.

The syntax $\pm x\pm y\le c$ is already the octagon domain
\cite[p.~1]{Mine}. Bozga--Iosif--Kone\v{c}n\'y's difference,
octagonal and finite-monoid affine relations have acceleration and
Presburger transitive-closure results
\cite[Definition 1 and Sections 4.2.1--4.2.2]{Bozga}. These principally
iterate a fixed relation and allow linear drift; here the relation
itself varies with $t$, and all previously output surface vertices
must remain distinct. A fixed-dimensional translation preserving that
condition has not been supplied. The term ``zigzag automata'' there
does not identify the subject with fullerene zigzags.

Jonoska--Krajcevski--McColm study counter languages of walks returning
to a unit cell in periodic digraphs \cite{Counter}; Wanke studies
path/cycle algorithms on finite periodic graphs \cite{Wanke}.
For these two sources (S2191, S2192), only the official abstracts and
bibliographic records were obtained. We do not infer their full
hypotheses or claim an application from an unread theorem. Walks or
decidability alone do not preserve literal simplicity or prove
periodicity of the parameter set.

\paragraph{Unary length is essential.}
The standard WS1S encoding of Presburger integers uses their binary
representations \cite[printed pp.~11--12]{Thomas}. Binary regularity
does not imply eventual periodicity of the represented integers:
the regular language $10^*$ represents powers of two. Our reduction
uses $n$ as unary word length, not the binary digits of $t$.
A separate formula for each $t$ would also be insufficient.

\paragraph{Boundary of the contribution.}
The received R005 comparison read the theorem statements and specified
sections listed in the packet's source records, not all proofs. Within
that reading, no prior result giving the full claim from the same
geometric hypotheses was located. Novelty is not established. The
remaining candidate contribution is the uniform, two-way translation
of geometric passages, seam ownership, legal truncation and all
surface-vertex collisions into one fixed MSO sentence, even with
growing widths and passage counts. Fixed-width theory and the known
logical conclusions do not supply that translation by themselves;
we do not prove that a shorter existing reduction is impossible.
The general convex-face proof is received research work, not an
already published theorem cited as external authority. Human proof
peer review and full comparison with catalogue/zigzag theory remain
necessary. The extension's unfinished geometric details do not
disappear on receipt of the convex-face proof.

The two-interval arithmetic kernel is covered by Li--Chen's uniform
bound for continued fractions of ratios of parameter polynomials
\cite[Theorem 27]{LiChen}. For several intervals, finite residue
branches and a uniform compressed node count, the stored survey did
not locate an identical general theorem; this is not a novelty proof.

K3's variable-modulus behaviour obstructs a single ordinary
multivariable Presburger description in fixed common coordinates.
Moving to canonical representatives or allowing variable divisibility
does not follow from raywise periodicity. The stronger M4/M4-prime
questions remain separate. General E10 $r=1$ necessity, untranscribed
families, all degenerations, human proof peer review, and a formal
geometric-to-program pipeline are not completed here. A3 minimum,
density and the IPR census belong to the first article. No new D12
classification is asserted by this paper.

\appendix
\section{Local polygons and zero boundaries}\label{app:domains}
\claimid{app:domains}
Twice total polygon area is the carbon order only when the quotient
is valid. The stored local coefficient audit gives
\begin{center}\small
\begin{tabular}{lll}\toprule
Signature & Side factors & Twice total area\\\midrule
E10 & $p,r,r+s$ & $12r^2+2s^2+12rs+24pr+12ps$\\
D12 & $s,r+s$ & $12r^2+48rs+48s^2$\\
D14 & $r,s$ & $8r^2+24rs+12s^2$\\
D1 & $p,r$ & $36p^2+48pr+12r^2$\\
K3 & $r,s$ & $14r^2+12rs$\\
B2 & $p,r$ & $12p^2+24pr+4r^2$\\
C1 & $a,b$ & $8a^2+16ab+4b^2$\\
O2 & $t,d$ & $14t^2+16td+2d^2$\\\bottomrule
\end{tabular}
\end{center}
The 71-chart/342-edge audit establishes local convexity, positive area
and side lengths in its stated positive domain, not global quotient
simplicity. With other parameters positive, E10 $s=0$ and D12 $r=0$
alone remain locally nondegenerate. E10 $p=0$ or $r=0$, K3 $r=0$,
B2 $p=0$ or $r=0$, C1 $a=0$, and O2 $t=0$ collapse some faces.
D12 $s=0$, either D14/D1 zero, K3 $s=0$, C1 $b=0$ and O2 $d=0$
produce zero sides. Removing an IET zero-length label is not the same
as redefining a graph after cone mergers. The omitted leapfrog-domain
data remain in the extension's prospective records, not the main
corollary. A3's $l=0$ and $k=l$ are the separate A1 and A2 boundaries.

\section{Representative formal statements and source correspondence}\label{app:correspondence}
The accompanying \record{v3-correspondence.md} maps every numbered v3
section to this reorganization, including the moved oblique-direction
sketch. It distinguishes translation, condensation and formal-status
updates; no prose-equivalence theorem is claimed.

The primary received declarations now include
\record{ClaudeChallenge2.main} and
\record{ClaudeChallenge2.State.main_statement} and
\record{ClaudeChallenge2.main_residue}, followed by
\record{Claude2.k3_ray_all}, \record{Claude2.e10_ray_all},
\record{Claude2.k3_example}, \record{Claude2.e10_example}, and
\record{Claude2K3.k3_sufficient}. The new K3 entries are
\record{ClaudeChallengeK3.main}, \record{Claude4K3.k3_statement} and
\record{Claude4K3.noThree}. The machine layer supplies their exact
types, source commit and the kind of proof receipt. The historical
representative kernel types remain:
\begin{quote}\small
\record{finite_literal_path_decision}: finite type, two-common-neighbour
hypothesis, cone predicate and fixed oriented triangle imply a first
repeat at most the number of vertices, and path existence iff a legal
endpoint occurs before it.\par
\record{cylinder_portal_shift}: a B\'ezout witness for primitive $u$
and $g\mid Q/\mathrm{cylHeight}(u,v)$ imply that the stated shifted
displacement is a multiple of $gu$. The integer-division and
divisibility hypotheses are retained.\par
\record{induction_uniform_depth}: a finite-list child relation with a
strictly decreasing lexicographic natural rank has a finite uniform
depth from each starting state.\par
\record{iet_component_uniform_depth}: the concrete equal-final-label
removal component has such a depth; the omitted IET branches are not
thereby formalized.
\end{quote}
The supplied exact sources, not these paraphrases, determine the Lean
types. Dependencies are pinned; caches, binaries, raw graph streams,
contact details are supplied separately from the proof sources.

\section*{Acknowledgements and AI-use statement}
The author used Codex and Claude Code
to assist with literature organization, mathematical exploration,
proof drafting, implementation, finite verification, formalization of
the explicitly identified arithmetic and graph statements and the
specification-level Theorem 1a, and English exposition. Claude workers
supplied the general theorem's and K3's Lean proofs against the
coordinator's prior specifications. The research team checked the
received sources and records; it did not supply an independent Lean
derivation or rerun Lean for this K3 edition. The author is responsible
for the published account, its provenance, and the stated limits of
verification. AI systems are not listed as authors. AI prose reviews
are identified separately from human review. Only the explicitly
listed Lean propositions are described as formally proved; finite
checks and independently constructed tests have their stated domains.
No claim of complete human checking of every argument is made.

\section*{License and contact}
Copyright 2026 Tatsuhiko Takimoto. The article text, figures and research
data are available under CC BY 4.0; the accompanying code and Lean
sources are under Apache-2.0, with third-party notices retained.
Contact is through GitHub Issues of the public repository; its URL
has one canonical location in the accompanying
\record{REPOSITORY-URL.txt}. No contact email is listed.

\clearpage
\begin{thebibliography}{99}
\bibitem{A3} Tatsuhiko Takimoto,
\emph{A two-stage classification of linear generalized Stone--Wales paths in
Goldberg--Coxeter images of the dodecahedron}, version 8.1 (2026).
{\urlstyle{rm}\url{https://github.com/tkmt-lab/a3-gsw-classification}}.
\bibitem{BBS} A. Bille, V. Buchstaber, and E. Spodarev,
\emph{Some open mathematical problems on fullerenes}, arXiv:2410.19322,
2024; published 2025, doi:10.1021/acs.jcim.4c01997.
\bibitem{Graver} J. E. Graver,
\emph{Catalog of All Fullerenes with Ten or More Symmetries},
DIMACS Series 69 (2005), 167--188; author version with corrections.
\bibitem{Connectivity} K. Kutnar and D. Maru\v{s}i\v{c},
\emph{On Cyclic Edge-Connectivity of Fullerenes}, arXiv:math/0702511v1,
2007. The received use is the cited established cyclic-connectivity
statement, not a new reading of its original proof.
\bibitem{MSO} D. Mandrioli, D. Martinenghi, A. Morzenti,
M. Pradella, and M. Rossi,
\emph{Lecture Notes on Monadic First- and Second-Order Logic on Strings},
arXiv:2301.10519v1, 2023, Section 3.1, Theorem 3.2.
\bibitem{BGW} T. Bogart, J. Goodrick, and K. Woods,
\emph{Parametric Presburger Arithmetic: Logic, Combinatorics, and
Quasi-polynomial Behavior}, Discrete Analysis (2017), article 4,
doi:10.19086/da.1254, Theorem 1.15.
\bibitem{QE} A. Mansutti and M. R. Starchak,
\emph{One-Parametric Presburger Arithmetic Has Quantifier Elimination},
LIPIcs 345, MFCS 2025, article 72, doi:10.4230/LIPIcs.MFCS.2025.72.
The stored reading covers the stated main theorems, not the full proof.
\bibitem{LiChen} N. Li and S. Chen,
\emph{On the Ring of Integer-valued Quasi-polynomials},
arXiv:0706.4288v2 (2007; internal PDF date 2008), Theorem 27.
The stored R005 reading distinguishes its nonzero-denominator hypothesis
from eventual positivity and from the several-interval setting.
\bibitem{Courcelle} B. Courcelle and J. Engelfriet,
\emph{Graph Structure and Monadic Second-Order Logic}, Cambridge
University Press, 2012; author-hosted Chapter 1, Theorems 1.12, 1.21.
\bibitem{Behmaram} A. Behmaram and C. Boutillier,
\emph{Asymptotic enumeration of perfect matchings in m-barrel fullerene graphs},
arXiv:1710.05156, Section 3.
\bibitem{Fischer} E. Fischer and J. A. Makowsky,
\emph{On spectra of sentences of Monadic Second Order Logic with counting},
author draft, 22 May 2003; Journal of Symbolic Logic 69 (2004), 617--640.
\bibitem{Kotek} T. Kotek and J. A. Makowsky,
\emph{Recurrence Relations for Graph Polynomials on Bi-iterative
Families of Graphs}, arXiv:1309.4020v1, 2013.
\bibitem{Mine} A. Min\'e, \emph{The Octagon Abstract Domain},
arXiv:cs/0703084v1, 2007, p.~1 only in the stored reading.
\bibitem{Bozga} M. Bozga, R. Iosif, and F. Kone\v{c}n\'y,
\emph{Fast Acceleration of Ultimately Periodic Relations},
VERIMAG Technical Report TR-2010-3, 2010.
\bibitem{Counter} N. Jonoska, M. Krajcevski, and G. McColm,
\emph{Counter machines and crystallographic structures},
Natural Computing 15 (2016), 97--113,
doi:10.1007/s11047-015-9527-0. Official abstract only.
\bibitem{Wanke} E. Wanke, \emph{Paths and cycles in finite periodic graphs},
MFCS 1993, LNCS 711, 751--760,
doi:10.1007/3-540-57182-5\_66. Official abstract only.
\bibitem{Thomas} W. Thomas, \emph{Languages, Automata, and Logic},
Report 9607, May 1996; printed pp.~11--12 (PDF pp.~14--15).
\end{thebibliography}
\end{document}
